\subsection{Tenth Letter}

\begin{quotex}
\begin{itemize}
\item
Guenon explains more about impeccability 
\item
Guenon expands on the powers of sorcerers and that no one is safe from their spells 
\item
Guenon offers opinions on Thomas Palamidessi, Carl Jung, Meher Baba, the Golden Dawn, and Aleister Crowley 
\item
Kerenyi sided with the liberals in Hungary during WW II. That may explain his change of attitude towards Guenon 
\item
Evola translated Leon de Poncins' book on the occult war into Italian. Both Guenon and Evola believed in conspiracies, aka ``counter initiation", which may explain their mutual interest in Poncins. 
\end{itemize}
\end{quotex}

29 October 1949

Cairo, Egypt

I received your letter of 4 September about eight days ago; I wonder if you were able to reach \textbf{Eliade}, since I rushed to send you his address in Capri. A short time after I wrote you, I knew that he had already returned to Paris; it seems that he did not take a very long vacation.

[Discussions and complaints about publishing and editors omitted]

\textbf{Abdul-Hadi}'s article, which I spoke about, entitled \emph{Pages dedicated to Mercury}, was republished in ``Etudes Traditionnelles", but I realized that occurred only after the war, which explains the reason why you aren't familiar with him.

As for ``impeccability", it goes without saying that it must belong equally to all who have reached a certain spiritual level. But, excluding the case of the prophetic mission, the possession of such a state concerns only the person who received it, and no one else can speak of its merit or care about it, hence the absence of every explicit affirmation in this regard in orthodox doctrine.

If I learn something about \textbf{Leon de Poncins}\footnote{\url{https://en.wikipedia.org/wiki/L\%C3\%A9on_de_Poncins}}, I will let you know, but I still don't know what happened to him. I believed, I no longer know exactly why, that he had to retreat to Switzerland during the war, but I was not able to confirm that.

Regarding evil spells, there is a great difference between true sorcerers like those with whom he had to deal and simple occultists. The latter, notwithstanding all their pretenses, never reach any effective result. There have often been some of them who attempted to do something against me and, also like you, I never heard anything about it at all.

On the other hand, when you think that things of that type should not be able to strike those who have a true spiritual vocation (however I don't think that can be said to have been that the case of Leon de Poncins), it is also necessary to make a distinction: if you want to speak of the psychic and mental side, you are absolutely correct, but things are quite different from the corporeal point of view and anyone can always be struck in this regard. Furthermore, since it has been passed down that some sorcerers succeeded in sickening the Prophet himself, I don't even see who could boast of being secure from their attacks.

\textbf{Thomas Palamidessi}, whose writings you inserted into your letter, is again obviously another charlatan of the type of those who currently abound everywhere. But what is astonishing is that he again appropriated ideas found in your and other's books, to use them in a way that can only discredit them; in such conditions, the works that he publishes should not require much effort to write.

I heard about \textbf{Meher Baba} in the past and his vow of silence, which does not seem to prevent him from responding in other ways to the questions that are asked of him, but I did not know that he has reappeared in recent times. I do not know if he ever was associated with any regular initiatic organization, but it seems dubious to me because he is a Parsi. Nothing of the type seems to exist among the Parsis of India, who moreover have conserved only rather incomplete fragments of their tradition (I speak of the Parsis of India, because those of central Asia have quite other knowledge, even if they keep it hidden).

I am quite astonished about how much you tell me in regards to \textbf{Karoly Kerenyi}\footnote{\url{https://en.wikipedia.org/wiki/K\%C3\%A1roly_Ker\%C3\%A9nyi}}, because I remember that in the past he had spoken very favorably about me; it had to be in 1939 or 1940, and at that occasion he had sent me his book \emph{Religion in Antiquity}. On the other hand I reviewed it, but because of the suspension of ``Etudes Traditionnelle", I was able to publish it only after its resumption.

As for \textbf{Carl Jung}, his influence unfortunately is gaining ground everywhere, in France as in Italy and Switzerland, and he seems to me still more dangerous than Freud because of his pseudo-spiritual pretense. Recently I had to write an article about the deformations of the very idea of Tradition provoked by his theory of the ``collective unconscious".

The Golden Dawn\footnote{\url{https://en.wikipedia.org/wiki/Hermetic_Order_of_the_Golden_Dawn}} was a self-styled Hermetic organization that fundamentally did not seem to have a very serious character, because it was from its beginnings an authentic mystification. It is true that this could serve to conceal some rather suspect things. Internally, the principle role was developed by MacGregor and his wife (Bergson's sister). Only much later was Crowley introduced to it, as he also did in many other things. Even when it was not about rather insignificant pseudo-initiations (perhaps that was not at all the case for the Golden Dawn), his involvement introduced truly sinister influences into it, if only making of it something much more dangerous. The Golden Dawn has ceased to exist, following a misunderstanding among its members, but a part of them followed it up under the name of Stella Matutina\footnote{\url{https://en.wikipedia.org/wiki/Stella_Matutina}}.

To come back to \textbf{Aleister Crowley}\footnote{\url{https://gornahoor.net/?p=6789}}, what you told me reminds me of the story that turned up in 1931 (I believe at least that was the exact date): while he was in Portugal, he suddenly disappeared. They found his clothes on the border of the sea, something that made them believe he had drowned. But it was only a simulated death, since they were no longer concerned about him and did not try to find out where he had gone. Actually, he went to Berlin to play the role of secret adviser to Hitler who was then at his beginning. It is probably this that had given rise to certain tales about the Golden Dawn, but in reality it was only about Crowley, because it does not seem that a certain English colonel named Etherton, who was then his ``colleague", had ever had the least relationship with that organisation.

A little later, Crowley founded the Saturn Lodge in Germany; have you ever heard of it? There he called himself Master Therion, and his signature was \emph{to mega Therion} (the Great Beast), something that in Greek gives exactly the numeric value 666.



\flrightit{Posted on 2012-08-20 by Cologero }

\begin{center}* * *\end{center}

\begin{footnotesize}\begin{sffamily}



\texttt{Charlotte on 2012-08-20 at 14:16 said: }

I'm somewhat loathe to comment again as I'm genuinely perturbed and often exasperated – to the point of intense sighing and face-palming – by some of what I've seen and heard on this website, albeit some very interesting subjects are covered. However, I feel this is an important matter for the esoteric community and want to say that I'm glad reservations about AC have been raised. Throughout my life, even since before I became consciously Christian, I have been intensely wary of him. I and was once told by a (male) witch to not even physically touch one of AC's books because his works are so intensely charged with `black magic’. (We were in a bookshop together and saw a upside-down pentacle-adorned tome looming in sinister fashion from the bottom shelf). I do realise that he is a very complex character – and pioneering in his own right – but his identification with the number of the beast says it all really. I also agree about the damage that can be inflicted by a sorcerer given the right conditions to work in – it would be possible to see the precarious state of the world and human spirituality as having been partially brought about by the nefarious actions of determined black magicians(I could name names in this regard but don't want to draw the wrong sort of attention to anyone here). AC's involvement in the war is a curious subject. Paradoxically in this context, I actually heard that he taught Churchill the famous `V' sign in order to counterbalance the dark magic of Hitler's gestures, but I'm not sure if that's just a legend. Chances are he was an Agent ZigZag character who worked for both sides. Either way he appears not to have won much respect for his contribution. Other magus-type characters, such as Gurdjieff, couldn't stand him, but magus types are notorious for not getting along – not enough space for more than one mammoth ego in one room I should imagine.


\hfill

\texttt{zero is zero on 2012-08-20 at 16:56 said: }

All this stuff is cool but where is the article on game Cologero? I can't wait. Is it going to be called ``Game and Sex Magic as Authentic Initiatory Tradition"?


\hfill

\texttt{Cologero on 2012-08-20 at 17:30 said: }

Mr. Zero, you need to check the poll results, which were overwhelmingly opposed to posts on Game.

Nevertheless, since I am not much of a fan of the democratic method, I wrote a sample article anyway, which was rejected by our Board. They felt it would be unwise to reveal publicly such intimate details of female psychology. I may try to rework it to get it approved.

In the meantime, while you're waiting, I suggest taking matters into your own hands.


\hfill

\texttt{Cologero on 2012-08-20 at 17:35 said: }

Thank you for your continued interest and valuable insights, Charlotte. Yet, I have no idea why you would be perturbed or exasperated by someone else's opinions. Enjoy what you find interesting and ignore the rest.

Perhaps Crowley is as evil as you say, but real evil seldom announces itself so clearly. It typically hides behind bright smiles and chatter about ``love", much like a wolf in sheep's clothing.


\hfill

\texttt{Charlotte on 2012-08-20 at 18:28 said: }

Well I haven't tended to think Gornahoor is about opinions, I think it is a group of people presenting their views as the wholly representative of the `primordial tradition'. which I do not necessarily see in the same way. Certainly many things are presented far more forcefully than as `opinions'. but there are nonetheless things we can agree on, if not all things. 

I know what you mean about wolves in sheep's clothing, opposites and mirrors, etc, there are even those who say that entities appearing `white' (in terms of colour) in the spiritual dimensions, including dreams, are in fact luciferic and manifesting evil. Personally I do not agree with this (not that I've been given many instances where the issue cropped up at all), although I WILL admit that the most evil entity I have ever encountered was indeed white, although this may be because it was translucent (ie, of no colour at all). 

Recently, in the interests of objectivity, I have tried to listen to some fans of Crowley and I can see that he did impart some truth within his body of work so I am not meaning to be overly judgmental, simply conveying what my own impression has been. I also really enjoyed his biography by Lawrence Sutin I think it was, as at the very least the man had a fascinating life and certainly did a lot for gay rights! The `evil sorcerer' I was referring to however, is not him, it is someone who was leader of the Temple of Set in the US. I don't really want to talk much further about that, especially here where there is the potential for the dark energy to be attracted, I just thought it was worth clarifying in case people thought I meant AC. In my opinion this (the activities of the temple leader) goes right to the dark heart of dark occultism in the modern world, but getting one's head around it, apart from being subjectively horrific, is so unbelievably complicated it's probably best left aside unless absolutely necessary.

It's also interesting what is said about Jung, as in the past few weeks I've been given cause to have doubts about him for the first time, although more for accusations of anti-semitism than anything else. As I think you know, I am anything but anti-semitic, so even a hint of anti-semitism is like anathema to me. Anyway, possibly Jung doesn't deserve that criticism, I didn't look any further into it but I was surprised.


\hfill

\texttt{Charlotte on 2012-08-20 at 18:32 said: }

details on female psychology!!! you fellas may need some help with that as it is not your strongest point (male psychology isn't my strong point either by the way) talk about a way to start a holy war…. :-))


\hfill

\texttt{Avery Morrow on 2012-08-20 at 18:34 said: }

I agree about Crowley: he is a charlatan and dangerous only insofar as he will waste people's time with false initiation. The real demons Charlotte should be looking out for are those like Screwtape, the henchmen of the occult war who work on our souls, notwithstanding Guénon's somewhat weird claim that wizards are attacking him.


\hfill

\texttt{Charlotte on 2012-08-20 at 18:57 said: }

I'm not looking!!! trying not to anyway. What led me to even try to learn something about AC was a decision I took to somehow resolve – as far as I might be able – the terrible situation so far as magic is concerned that was brought about during and around the World Wars. It seemed very clear to me that the present day spiritual dangers – and especially the impasse – would never be resolved until that particular `knot' was untied. To say I wasn't prepared for what I saw would be an understatement, Guenon isn't the only one who reckons he's had wizard's past, present and/or future attacking him! (although voodoo, of course, requires compliance on the part of the victim, which is why hypnotism is so widely employed as a dark art). What you say about the henchmen who work on our souls is very true, terrifyingly so, it's important to be protected, the armour of God is necessary. I also should say that certain spiritual states of being make one more vulnerable to `attack' from `wizards'. with relative degrees of manifestation that can come right through to the physically injurious. In very heightened states thought becomes so much more powerful – what might otherwise be `undesirable' but `harmless' thoughts – anger with another, or lust, for example – take on a terrifying potential when you're in astral/etheric realms. If the `victim' and `attacker' are both – possibly inadvertently – at a very high stage of spiritual development or psychic awareness then you can quadruple it all again or thereabouts. look into the abyss and it looks back etc….

This is also the reason why `wizard' or would-be-wizard types, for want of a better term, speak in a certain way about their exploits and or experiences, as if they were as tangibly real in a scientific and objective sense as someone eating an apple. T


\hfill

\texttt{Charlotte on 2012-08-20 at 19:05 said: }

To give a more concrete example of what Guenon may be referring to: I know a woman who had a relationship with a Muslim man a few years ago. Unbeknownst to my friend, it turns out there was a Muslim woman vying for his attention, who went to an Imam in her home country and confided to him her problem. If I remember rightly she was given some sort of bracelet, something with three beads, each of which represented a misfortune that was to befall her love rival. Sure enough, my friend duly lost around half her body weight, almost all of her hair and fell into a deep depression that was only alleviated when she converted to Islam (via a different man, who would become her husband) and found her own Imam to lift the spell…


\hfill

\texttt{zero is zero on 2012-08-20 at 19:06 said: }

Haha good answer Cologero.

Re Crowley it's very obvious that something went wrong somewhere along the road with him, someone who bases his whole philosophy around Will but dies as a heroin addict, I think anyone can see there is a major problem involved. That doesn't mean that nothing worthwhile can be gained from some of his writings. My feeling is that we are doomed to find knowledge in this way, making a collage from scatttered sources, finding kernels of truth hidden here and there, there is no unified body of knowledge to act as a whole source in sight.


\hfill

\texttt{Charlotte on 2012-08-20 at 19:24 said: }

The only way is by revelation from the Holy Spirit and it's probably best if one can be satisfied with that gift rather than trying to understand from an intellectual point of view. To take the life rather than the knowledge, as it were. However human beings are curious by nature and it appears to be almost impossible for almost anyone to stay firm in this respect (the trouble being that the mind is where doubt can be instilled).


\hfill

\texttt{Charlotte on 2012-08-20 at 19:33 said: }

By the way I DO think evil tends to appear as it is – just that, Evil – but the trouble is knowing how, when and why it appears. I have never heard evil chattering about love, although I do know someone who often talks about love to mask her own despair. Nobody need be fooled by smiles – in Thailand, for instance, if you get into an argument with someone and they start smiling then watch out, they may be about to kill you.


\hfill

\texttt{Cologero on 2012-08-20 at 22:53 said: }

Charlotte, with all the time you have allegedly spent reading Gornahoor, you still think it is all about the ``mind". Shame on you.


\hfill

\texttt{Cologero on 2012-08-20 at 23:01 said: }

Au contraire, ma chère Charlotte, I do understand female psychology, inasmuch as I knew you would come back.


\hfill

\texttt{Charlotte on 2012-08-21 at 05:02 said: }

I had to give you the pleasure of sticking the knife in one more time, eh! pride thankfully isn't my main vice…you might have been able to judge my reaction, seeing as you know how much I care, but you will never be able to understand the reasons why. Enjoy the Miami sun, it's raining here in England but at least we've not dried out!


\hfill

\texttt{Charlotte on 2012-08-21 at 05:37 said: }

By the way, where did I say I thought Gornahoor was `all about the mind'. it was YOU who suggested it was a mere collection of opinions, I actually gave it more credence than that, but perhaps I shouldn't have? 

While I am here, I may as well say I couldn't help noticing that people have been confused about `Guenon's Western Initiation' and which `school' it might have been. Well, he may have joined some other organisation in England or wherever – Druidism, perhaps – but at least I can tell you that Sufism IS a fundamental aspect of the western mystery tradition, along with Christianity and Judaism. It's easiest to see them as three interlocking circles if one would like a symbol representing how they fit together. 

The `Eastern' tradition is along the Buddhist lines. The fundamental difference is that the Western mysteries are revealed from the Holy Spirit and keep with the idea of a Creator God that is distinct from the creation, whereas the Eastern view sees Creator and Created as being one, without distinction. 

This is a particularly important difference when you start looking at concepts of eternity, of which there are more than one.


\hfill

\texttt{Michael on 2012-08-21 at 11:35 said: }

By the way Cologero, I second Zero's request for Game posts. It seems that most of the commenters are far along reading works of Tradition, but more practical posts will help novices, such as myself, come up to speed faster.


\hfill

\texttt{Cologero on 2012-08-21 at 12:45 said: }

Michael, the Board has approved my posts conditionally, provided I file them under ``fiction". They may not be what you expect.

The Board also suggested we wait until Charlotte goes away again, lest things get out of hand.


\hfill

\texttt{Cologero on 2012-08-21 at 13:03 said: }

No, Charlotte, there was a ``West" long before Christianity or Islam came on the scene, and Judaism–whatever that is–had no influence on it. Repeating some nonsense from a Karen Armstrong book is not going to cut it here.

Does ``the Eastern view sees Creator and Created as being one"? Perhaps in the sense that Nirvana and Samsara are one? How about the Thomist view that for God, Essence and Existence are one?

Either all is one and there are no important differences, or all is not one, which makes one true and the other false.

As for your other point, there is only one concept of eternity, but many ways to misunderstand it.

I am asking you nicely, Charlotte, to focus on what is said here. If anything gets you emotionally worked up, that is because you are experiencing resistance, more precisely, resistance to gnosis. The first task for any reader here is to clear his mind of all personal opinion in order to hear something new.

That is also the requirement of the second Arcana: the mind must be freed of all perturbations, specifically including strong negative emotional reactions, as well as the putrefaction of useless opinions. Only then can the Holy Spirit be heard.


\hfill

\texttt{Charlotte on 2012-08-21 at 15:29 said: }

I haven't read a Karen Armstrong book, but you clearly have – what does she write about? 

I did focus on what was said here, which was Guenon's point about Crowley. It was YOU who chose to make inaccurate assessments of comments I'd written by wrongly saying I thought Gornahoor was all about the mind – again your assumptions, not my words. You can be as unpleasant towards me as you like – nothing new there – but it won't make what you say true. Your personal dislike of me is clouding your judgment. Far from being over emotional, I was at least capable of rising above my own reservations by agreeing with something that was presented here when I felt it was important enough to merit comment. 

The western mystery tradition, having evolved over many years, IS informed by Sufism, Judaism and Christianity. Like it or not. 

Go ahead and print your article about women, none are likely to dare comment so yourself and the `Board' need never by any the wiser. Why don't you become a closed group if you wish to avoid attracting unwanted attentions from those you see as lesser mortals?


\hfill

\texttt{Charlotte on 2012-08-21 at 16:19 said: }

Seeing as you wrongly assumed I'd been quoting an author I barely know, here is some pertinent information from Gareth Knight (a living male authority on the western mystery tradition, no less!). He talks more about the difference between West and East so far as mysteries are concerned, it's been up on my blog for ages:

The pendulum has swung back again—or at any rate is about to start its swing. In speaking of `the unity of the world and all things in it'. we must, however, avoid the error of oriental monism which denies the dual existence of Creator and created. According to this view the universe and all the inner worlds therein have been self-created, or at best emanated from a central source.

This means that God is in everything, in the holiest of holies and in the dust on the sandals of the worshipper at the temple gate. As a child of an acquaintance put it with devastating childlike logic. `When I stamp on the ground am I stamping on God?' To this the monist would rush to reply `Yes'. but the theist would say `No'. The monist would go on to say that as God is also in the child's foot, sock and shoe, God was stamping on God. The theist would go on to say that although God is not in everything He is omniscient as far as the creation is concerned and is therefore aware of the child stamping and in empathy with both the child and the ground.

All this is not academic, theological or philosophical hair splitting, for the consequences of believing one thing or the other are profound. If we are going to build a philosophical or theological edifice we need to be very certain of the rock upon which it is founded. To believe that all things unfurl of their own accord from nothing is to assume that man is capable of expanding his consciousness until he comes eventually as God, comprehending all — and that animals expand their consciousness to become humans, plants likewise to become animals, even minerals to become plants.

This is a theory that is, in fact, held by many students of the occult, based on the monist philosophical assumptions of the East It has its superficial attraction as a logical sounding kind of arrangement. It takes in the ideas of human progress and general life evolution that were newly formulated and current in the nineteenth century, and it is hardly surprising that these ideas in occult form were first promulgated in the West in the late nineteenth century by the efforts of the newly formed Theosophical Society.

What Madame Blavatsky, its founder, did really was to take nineteenth-century materialist evolutionary theory as formulated by Darwin and stand it on its head as a spiritual evolutionary theory, in much the same way that

Marx had inverted the spiritual dialectic of Hegel to form the dialectical materialism of Marxism. Both Marxism and Theosophy have a great spurious appeal as seeming to answer many questions by this agile topsy-turveydom. Unfortunately both are wrong — though this does not alter the fact that Marxism as a political philosophy came to dominate a third of the world and Theosophical monism dominates much of modern occult thought.

It is not our task to try to judge why certain particular nineteenth-century philosophical ideas should retain such a hold into modern times, though in the case of oriental monism and occultism its influence spread because a whole generation of occult students sat at the feet of Madame Blavatsky and imbibed her principles even if they later rejected some of the superstructure of her philosophy. They later taught others and so the basic assumptions spread — with various modifications to and arguments about the superstructure, but with the entire theological foundations taken for granted and accepted unchallenged.

The whole Western occult tradition, which had followed an underground course for centuries, burst out into the open, only to be thoroughly mixed, swamped and diluted with Eastern ideas deriving from Hinduism and Buddhism. The true occult heritage of the West stems, however, along with the religion of the West, from Christian and Judaic tradition — or rather from revealed as opposed to natural religion.

Gareth Knight, Experience of the Inner Worlds, The Sphere of Light


\hfill

\texttt{Charlotte on 2012-08-21 at 16:44 said: }

Now I will quote you, Cologero:

``Judaism – whatever that is"

Shame on YOU. Avery was right to caution about Screwtapes, casual readers beware.


\hfill

\texttt{Cologero on 2012-08-22 at 07:34 said: }

Charlotte, Judaism has a long history. At one time there were priests, temples and animal sacrifice. That judaism no longer exists. Rabbinic judaism is different. Orthodox Jews are different from Reformed Jews. And so on.

The same for Christianity … that word means very little today as a guide to what someone believes or how he acts.

You persist, Dear Charlotte, in finding the worse possible interpretation of what you read. St. Ignatius Loyola recommended the opposite, i.e., first consider the best interpretation of what you read before you judge.

It is time for you to leave us in peace, Charlotte, your sort of cattiness is not welcome here.


\hfill

\texttt{harlotte on 2012-08-22 at 07:47 said: }

I'm afraid you have again misjudged me Charles, the cattiness and unpleasant comments came from you. The actual decision to repost here, apart from the interesting comment about black sorcerers, was following a family funeral where a conversation with my catholic cousin gave me hope for the first time in months that not all is lost for the church. My tentative venture here proved once again how highly unpleasant and totally unreasonable professed `traditionalists' can be. I'm sorry you can't handle the truth I bring you.


\hfill

\texttt{Cologero on 2012-08-22 at 08:00 said: }

Thanks for clearing that up, Charlotte, and I am sorry for your loss.

Should you ever decide to comment on what is actually written here, you will be most welcome.


\hfill

\texttt{Charlotte on 2012-08-22 at 08:19 said: }

That's alright,you're not expected to know the ins and outs of my life. But please see that I WAS actually commenting on what is written here (Guenon on Crowley), not just turning up to be difficult – you can't seriously imagine I would subject myself to running the gauntlet on a whim, for the sake of being shot down! 

The truth is I thought this was something that we were able to agree upon and therefore a `safe topic'. In my view it was worth highlighting as it touches on matters pertinent to the `occult war'. I know that email/internet makes it very difficult to judge people's tone or the manner in which comments are intended, but please be assured that cattiness and sarcasm are simply not my style, although I can't help being ironic at times…i guess what is blatantly obvious to one person (any person), is shrouded in obscurity for another, we all see what we see.

As for `western initiations'. I will concede that it is fair to say there are a number of takes on what this entails. As far as I'm concerned it encompasses the revealed religion of Judaism, Christianity and Islam, which as you point out have all been through various changes and interpretations over the centuries. However, others see the `western mysteries' as a more recent phenomena, a Golden Dawn term, or see it primarily in terms of Egypt, for example. So there is potential for confusions, certainly.

It just seemed to me that so much is made of Guenon being Sufi that it was logical to assume this constituted his `western initiation'. but maybe there was another. Given his feelings on Crowley it seems unlikely it was Golden Dawn or such like, but other may know better. There were plenty of occult societies in those days so I guess it could have been anything…


\hfill

\texttt{Andrew on 2012-08-23 at 06:26 said: }

I honestly think we should just ban women. This is typical, and it's not Charlotte's fault, but…she is a woman, and we can't expect any better.


\hfill

\texttt{Cologero on 2012-08-23 at 08:27 said: }

On the other hand, Andrew, there are Buddhist texts, at least in the Vajrayana scriptures, in which women appear as teachers. So I suggest we take this opportunity to recognize and eradicate ignorance, irrationality, and unjustified aggression within ourselves.


\hfill

\texttt{Aghorable on 2012-08-23 at 10:02 said: }

Were I chief of a bricks-and-mortar club, I would perhaps adopt Andrew's suggestion, tempering it slightly to allow women with special skills to join us on certain nights.

Skills, for instance, belly dancing, harp playing, tea brewing, rhythmic gymnastics etc. They would only be permitted to indulge us in romantic discourse, no crusty erudition or intellectual posturing, but they would need to express a basic familiarity with metaphysical concepts to gain admittance, in order to appreciate the elevated states of our being.

And there would be hashish and opium and wine, and the finest blends of tobacco known to man, and sturdy eunuchs from barbarous corners of the world tending to our vagaries, with water flowing perpetually from a hidden mountain spring, and soft cushions of fine silk sent us by the Thai royal family; juicy fruit and tropical birds tamed into submission by our otherworldly song, and unlimited broadband internet to facilitate our studies, powered by thorium reactors…

Ah, Captain Nemo, my old friend.


\hfill

\texttt{Charlotte on 2012-08-23 at 11:02 said: }

Oh I forgot to mention the Rishis and Dakinis, that was somewhat remiss of me! 

Gosh what a dull world it would be without them…


\hfill

\texttt{Matt on 2012-08-23 at 20:00 said: }

….Yikes. To bring this back to Guenon's letter (the original topic at hand for this discussion), Guenon's judgement of Palamidessi is rather noteworthy…if only because Palamidessi's writings were the subject of a couple of posts on gornahoor by Logres. If I remember correctly Logres was planning on writing some more posts on this particular figure's writings, though they never materialized.

I wonder if Evola may have eventually come to share Guenon's view on Thomas, though I think Evola wrote a relatively cordial letter to Palamidessi late in his life…..so maybe not.


\hfill

\texttt{Logres on 2013-07-14 at 10:22 said: }

I was doing some rethinking the other day, and realized that Palamidessi's conversion to Christianity was in the 1960s: this letter is from substantially earlier. Is it possible that Tomasso was a charlatan, but one who converted and repented? Not necessarily so, just because he was formally a Christian later in life, but a possibility.


\hfill

\texttt{Sy on 2017-06-23 at 08:02 said: }

the ``article about the deformations of the very idea of Tradition provoked by his theory of the collective unconscious" would be interesting. has anyone found it?


\hfill

\texttt{Cologero on 2017-06-23 at 14:53 said: }

@Sy, see Chapter 5, ``Tradition and the Unconscious" in Guenon's \emph{Symbols of Sacred Science}.


\end{sffamily}\end{footnotesize}
